
Semantic Role Labelling~\citep[SRL, ][]{marquez08srl} is generally understood as 
the task of identifying and classifying the semantic arguments and modifiers of 
the predicates mentioned in a sentence. For example, in the case of the 
following sentence\footnote{``Haag plays Elianti'' is a segment of a sentence in 
training corpus.}:
\begin{quote}
\begin{center}
    \includegraphics[scale=.63]{haag-example}
\end{center}
\end{quote}
we are to find out that for the predicate token {}``plays'' with sense ``play a 
role'' (play.02) the phrase headed by the token {}``Haag'' is referring to the 
player (A0) of the play event, and the phrase headed by the token {}``Elianti''  
is referring to the role (A1) being played. SRL is considered as a key task for 
applications that require to answer {}``Who'', {}``What'', {}``Where'', etc.  
questions, such as Information Extraction, Question Answering and Summarization.  

In this paper we present a Markov Logic \citep[ML,][]{richardson06mln} to define 
a multi-language joint SRL model. We present a brief introduction to ML in 
section \ref{sec:markovlogic}.  At the core of ML there is a set of First Order 
Logic (FOL) formulae which models the SRL task. And advantage of the ML logic 
framework is that the FOL formulae can be define the in a modular fashion.  With 
this in mind we define a modular ML model which captures the different stages of 
the task: argument identification, argument classification, and sense 
disambiguation. We also define modules for some of language specific aspects of 
the task. We present the types and modules of our ML model in section 
\ref{sec:model}. These ML modules are the building blocks that we use to create 
a ML model for a specific language. Once we put together a ML model for a 
specific language we can develop a joint probabilistic model over all decisions 
involved int the SRL task. We present our experiments and result for each of the 
languages of the task in section \ref{sec:results} and in section 
\ref{sec:analysis} we present some error analysis.

